\subsection{定義と性質}

$R$をcommutative ring, $M$を$R$-moduleとする．
さらに，$\iota$を(有限とは限らない)添え字集合，$N$を$R$-moduleとする．

\begin{definition}
  \label{def:RootPairing}
  \lean{RootPairing}
  \leanok
  $M, N$は$R$上のperfect pairingとする．
  すなわち，bilinear map $B : \map{M \times N}{R}$が与えられていて，片方の変数を固定して単にlinear mapと見たとき，それぞれがbijectiveであるとする．
  このとき，$\Phi := \set{\alpha_i \in M}{i \in \iota},\ \Phi^\vee := \set{\alpha^\vee_i \in N}{i \in \iota}$が次を満たすとき，$P := (\Phi, \Phi^\vee)$を\textbf{root pairing}という：
  \begin{enumerate}[label=(\roman*)]
    \item $\forall i,\ B(\alpha_i, \alpha^\vee_i) = 2$ \label{defi:RootPairing_EqTwo}
    \item $\forall i j,\ \exists k,\ \alpha_j - B(\alpha_j, \alpha^\vee_i) \alpha_i = \alpha_k$
    \item $\forall i j,\ \exists k,\ \alpha^\vee_j - B(\alpha_i, \alpha^\vee_j) \alpha^\vee_i = \alpha^\vee_k$
  \end{enumerate}
\end{definition}

\begin{definition}
  \label{def:RootPairing.IsCrystallographic}
  \lean{RootPairing.IsCrystallographic}
  \leanok
  $P$をroot pairingとする．
  $\forall i, j,\ B(\alpha_i, \alpha^\vee_j) \in \Z$のとき，$P$は\textbf{crystallographic}であるという．
\end{definition}

\begin{definition}
  \label{def:RootPairing.IsIrreducible}
  \lean{RootPairing.IsIrreducible}
  \leanok
  $R$-加群$M, N$に関するroot pairing $P$が次を満たすとき，$P$は\textbf{irreducible}であるという：
  \begin{enumerate}[label=(\roman*)]
    \item $M$は非自明である，すなわち$M \ne 0$.
    \item $N$は非自明である，すなわち$N \ne 0$.
    \item $M$の部分加群$q$が，任意のreflectionで不変(i.e., $s_i(q) \subseteq q$)で，$q \ne 0$であれば，$q = M$である．
    \item $N$の部分加群$q$が，任意のcoreflectionで不変で，$q \ne 0$であれば，$q = N$である．
  \end{enumerate}
\end{definition}